\definecolor{dsdlcodebg}{RGB}{232,238,245}
\definecolor{dsdlcodeborder}{RGB}{208,215,222}
\newsavebox{\dsdlcodebox}
\newenvironment{dsdlcodeblock}
{%
  \par\noindent
  \begin{latin}
  \begin{lrbox}{\dsdlcodebox}
  \begin{minipage}{1.05\linewidth}
  \small
}
{%
  \end{minipage}
  \end{lrbox}
  \setlength{\fboxsep}{6pt}
  \noindent\fcolorbox{dsdlcodeborder}{dsdlcodebg}{\usebox{\dsdlcodebox}}
  \end{latin}
  \par
}

\section{مقدمه}

در حافظه‌ی معمولی، یک آدرس برای خواندن محتوای محل مشخصی ارائه می‌شود. در حافظه‌ی
انجمنی\LTRfootnote{Content Addressable Memory}، جهت جست‌وجو برعکس است: داده
جست‌وجو می‌شود و حافظه مشخص می‌کند کدام ورودی‌ها با آن تطبیق دارند. این ویژگی برای
مسائلی مناسب است که تصمیم باید بر اساس محتوای کلید و نه آدرس از پیش معلوم انجام شود.

\لر{TCAM}\LTRfootnote{Ternary Content Addressable Memory} نوع توسعه‌یافته‌ی
\لر{CAM} است. در \لر{CAM} معمولی هر بیت فقط صفر یا یک است، اما در \لر{TCAM} حالت
سوم \لر{X} نیز وجود دارد. \لر{X} به معنای بی‌تفاوت است و بیت متناظر کلید جست‌وجو را
از مقایسه حذف می‌کند. صورت آزمایش طراحی یک \لر{TCAM} شامل ۱۶ ثبات ۱۶‌بیتی را می‌خواهد
و تأکید می‌کند که طرح باید پارامتری باشد.

\section{تعریف تطبیق سه‌حالته}

برای هر بیت یک ورودی ذخیره‌شده، دو مقدار نگه‌داری می‌شود:

\شروع{فقرات}
\فقره \کد{stored\_data} مقدار صفر یا یک پایه را نگه می‌دارد؛
\فقره \کد{stored\_x\_mask} مشخص می‌کند کدام موقعیت‌ها \لر{X} هستند؛
\فقره اگر بیت ماسک صفر باشد، بیت داده باید با \کد{search\_data} برابر باشد؛
\فقره اگر بیت ماسک یک باشد، نتیجه‌ی آن موقعیت مستقل از \کد{search\_data} برابر تلقی می‌شود.
\پایان{فقرات}

برای یک ورودی با عرض $W$، بردار اختلاف مؤثر از رابطه‌ی زیر به‌دست می‌آید:

\begin{equation}
\mbox{\کد{mismatched\_bits}} =
(\mbox{\کد{search\_data}} \oplus \mbox{\کد{stored\_data}})
\land \neg \mbox{\کد{stored\_x\_mask}}
\end{equation}

ورودی زمانی تطبیق دارد که معتبر باشد و هیچ اختلاف مؤثری باقی نماند:

\begin{equation}
\mbox{\کد{match}} = \mbox{\کد{valid}} \land
(\mbox{\کد{mismatched\_bits}}=0)
\end{equation}

این نمایش دو مزیت دارد. نخست آنکه \لر{X} به‌عنوان مقدار ناشناخته‌ی شبیه‌ساز وارد مدار
سنتزپذیر نمی‌شود. دوم آنکه قاعده‌ی مقایسه با چند عمل منطقی ساده قابل پیاده‌سازی است.

\section{مشخصات و رابط طرح}

هر دو پیمانه‌ی \کد{tcam} و \کد{tcam\_no\_loop} دو پارامتر دارند:

\شروع{فقرات}
\فقره \کد{DATA\_WIDTH}: عرض هر ورودی، با مقدار پیش‌فرض ۱۶؛
\فقره \لر{DEPTH}: تعداد ورودی‌ها، با مقدار پیش‌فرض ۱۶.
\پایان{فقرات}

جدول~\رجوع{جدول:رابط} نقش سیگنال‌های خارجی را خلاصه می‌کند.

\شروع{لوح}[ht]
\تنظیم‌ازوسط
\شرح{رابط خارجی پیمانه‌ی \لر{TCAM}}
\شروع{جدول}{|c|c|p{7.6cm}|}
\خط‌پر
\سیاه سیگنال & \سیاه جهت & \سیاه نقش \\
\خط‌پر \خط‌پر
\لر{clk} & ورودی & ساعت ثبت داده و ماسک \\
\لر{reset} & ورودی & بازنشانی فعال‌بالا و هم‌زمان برای بیت‌های اعتبار \\
\کد{write\_enable} & ورودی & بردار \لر{DEPTH}‌بیتی؛ مجوز مستقل هر ورودی \\
\کد{write\_data} & ورودی & بردار داده با عرض \کد{DEPTH} ضرب در \کد{DATA\_WIDTH} \\
\کد{write\_x\_mask} & ورودی & بردار ماسک با همان عرض؛ یک بخش مستقل برای هر ورودی \\
\کد{search\_data} & ورودی & کلید مشترک جست‌وجو با عرض \کد{DATA\_WIDTH} \\
\کد{match\_lines} & خروجی & یک بیت نتیجه برای هر ورودی \\
\خط‌پر
\پایان{جدول}
\برچسب{جدول:رابط}
\پایان{لوح}

در پیکربندی پیش‌فرض، \کد{match\_lines} شانزده‌بیتی است. بیت $i$ این بردار مستقیماً
نتیجه‌ی ورودی $i$ را نشان می‌دهد. اولویت‌گذاری میان تطبیق‌ها انجام نمی‌شود؛ بنابراین چند بیت
می‌توانند هم‌زمان یک باشند. این رفتار برای مثال صورت آزمایش ضروری است، زیرا یک داده باید با
سه الگوی متفاوت تطبیق پیدا کند.

\section{معماری کلی}

طرح از یک پیمانه‌ی ذخیره‌سازی و دو پیاده‌سازی سطح بالا تشکیل شده است:

\شروع{فقرات}
\فقره \کد{tcam\_entry.v}: ذخیره‌ی داده، ماسک و بیت اعتبار یک ورودی و تولید نتیجه‌ی همان ورودی؛
\فقره \کد{tcam.v}: ساخت تعداد پارامتری از \کد{tcam\_entry}ها با حلقه‌ی \لر{generate}؛
\فقره \کد{tcam\_no\_loop.v}: اتصال همان ورودی‌ها با آرایه‌ی نمونه‌ها و بدون حلقه‌ی \کد{for}.
\پایان{فقرات}

هر ورودی، بیت مجوز و بخش داده و ماسک مستقل خود را دریافت می‌کند. بنابراین چند ورودی
می‌توانند در یک لبه‌ی ساعت، مقادیر متفاوت بگیرند. \کد{search\_data} به همه‌ی ورودی‌ها
متصل است و مقایسه‌ها به‌صورت موازی انجام می‌شوند.

نوشتن ترتیبی و جست‌وجو ترکیبی انتخاب شده است. داده و ماسک فقط در لبه‌ی بالارونده‌ی ساعت
تغییر می‌کنند، اما پس از تغییر \کد{search\_data}، نتیجه‌ی تطبیق بدون انتظار برای ساعت جدید
به‌روزرسانی می‌شود. این تصمیم با ماهیت حافظه‌ی انجمنی سازگار است و زمان پاسخ جست‌وجو را به
تأخیر مسیر مقایسه محدود می‌کند.

\section{طراحی یک ورودی \لر{TCAM}}

هر \کد{tcam\_entry} سه ثبات \کد{stored\_data}، \کد{stored\_x\_mask} و
\لر{valid} دارد. هنگام \لر{reset} فقط \لر{valid} صفر می‌شود. پاک‌نکردن فیزیکی داده و
ماسک مشکلی ایجاد نمی‌کند، زیرا \لر{valid} صفر، خروجی \لر{match} را بدون توجه به محتوای
قدیمی غیرفعال نگه می‌دارد. این تصمیم منطق بازنشانی را کاهش می‌دهد و از پاک‌سازی غیرضروری
تمام بیت‌های ذخیره جلوگیری می‌کند.

خطوط ۲۳ تا ۳۵ پرونده‌ی \کد{tcam\_entry.v} منطق نوشتن، ماسک‌کردن و تطبیق را نشان می‌دهند.
شماره‌های کنار کد همان شماره‌خط‌های پرونده‌ی اصلی هستند.

\begin{dsdlcodeblock}
\begin{verbatim}
23  always @(posedge clk) begin
24      if (reset) begin
25          valid <= 1'b0;
26      end else if (write_enable) begin
27          stored_data <= write_data;
28          stored_x_mask <= write_x_mask;
29          valid <= 1'b1;
30      end
31  end
32
33  // Masked positions are removed before checking for unequal bits.
34  assign mismatched_bits = (search_data ^ stored_data) & ~stored_x_mask;
35  assign match = valid && (mismatched_bits == {DATA_WIDTH{1'b0}});
\end{verbatim}
\end{dsdlcodeblock}

هنگام فعال‌بودن \کد{write\_enable} در لبه‌ی بالارونده، داده و ماسک جدید ذخیره و ورودی
معتبر می‌شود. سپس منطق ترکیبی اختلاف بیت‌ها را محاسبه می‌کند. ابتدا \لر{XOR} تفاوت
\کد{search\_data} و \کد{stored\_data} را مشخص می‌کند؛ سپس معکوس ماسک فقط بیت‌های مهم
را نگه می‌دارد. صفرشدن کل بردار اختلاف به معنای تطبیق کامل همه‌ی موقعیت‌های غیر \لر{X} است.

\section{طراحی پیمانه‌ی سطح بالا}

هر دو پیمانه از \کد{tcam\_entry} استفاده می‌کنند. هر بیت فعال \کد{write\_enable}،
نوشتن ورودی متناظر را مجاز می‌کند. اگر بازنشانی فعال نباشد، ورودی‌های غیرفعال داده،
ماسک و اعتبار قبلی خود را حفظ می‌کنند.

عرض هرکدام از \کد{write\_data} و \کد{write\_x\_mask} برابر حاصل‌ضرب تعداد ورودی‌ها در عرض هر ورودی است. برش ورودی با شماره $i$ از بیت حاصل‌ضرب $i$ در عرض داده آغاز می‌شود؛ ورودی صفر در کم‌ارزش‌ترین بخش قرار می‌گیرد. برای چهار ثبات سه‌بیتی، اتصال \کد{\{3'b111, 3'b100, 3'b010, 3'b001\}} به \کد{write\_data}، به‌ترتیب داده‌های \کد{001}، \کد{010}، \کد{100} و \کد{111} را به ورودی‌های صفر تا سه می‌رساند. با مجوز \کد{1111} هر چهار مقدار در همان لبه ثبت می‌شوند؛ با مجوز \کد{1010} فقط ورودی‌های یک و سه تغییر می‌کنند. کلید \کد{search\_data} همچنان مشترک است و همه تطبیق‌ها در \کد{match\_lines} دیده می‌شوند.

در \کد{tcam.v} حلقه ساخت، برش داده، برش ماسک و بیت مجوز هر ورودی را به نمونه مربوط وصل می‌کند. در \کد{tcam\_no\_loop.v} آرایه نمونه‌ها همین اتصال را بدون \کد{for} می‌سازد. بردارهای داده و ماسک میان نمونه‌ها تقسیم می‌شوند، ولی ساعت، بازنشانی و کلید جست‌وجو مشترک می‌مانند. این دو شیوه رفتار خارجی یکسان دارند؛ حلقه ساخت نیز مراحل زمانی پیاپی ایجاد نمی‌کند.

اتصال‌های نسخه حلقه‌دار در خطوط ۲۵ تا ۳۵ پرونده \کد{tcam.v} چنین است:

\begin{dsdlcodeblock}
\begin{verbatim}
25  for (row = 0; row < DEPTH; row = row + 1) begin : entries
26      tcam_entry #(.DATA_WIDTH(DATA_WIDTH)) entry (
27          .clk(clk),
28          .reset(reset),
29          .write_enable(write_enable[row]),
30          .write_data(write_data[row*DATA_WIDTH +: DATA_WIDTH]),
31          .write_x_mask(write_x_mask[row*DATA_WIDTH +: DATA_WIDTH]),
32          .search_data(search_data),
33          .match(match_lines[row])
34      );
35  end
\end{verbatim}
\end{dsdlcodeblock}

آرایه نمونه‌ها\LTRfootnote{Instance array} در خطوط ۲۲ تا ۳۰ پرونده
\کد{tcam\_no\_loop.v} همان اتصال را بدون حلقه بیان می‌کند:

\begin{dsdlcodeblock}
\begin{verbatim}
22  tcam_entry #(.DATA_WIDTH(DATA_WIDTH)) entries [DEPTH-1:0] (
23      .clk(clk),
24      .reset(reset),
25      .write_enable(write_enable),
26      .write_data(write_data),
27      .write_x_mask(write_x_mask),
28      .search_data(search_data),
29      .match(match_lines)
30  );
\end{verbatim}
\end{dsdlcodeblock}

\subsection{محدوده پارامترها و حالت صفر‌بیتی}

محدوده معتبر هر دو پیمانه، عرض داده و تعداد ورودی بزرگ‌تر یا مساوی یک است. پیکربندی یک ثبات صفر‌بیتی عمداً رد می‌شود. عبارت \کد{[-1:0]} در \کد{Verilog} بردار صفر‌بیتی نیست، بلکه یک بازه دو‌بیتی است؛ بنابراین تفسیر آن به‌عنوان حافظه صفر‌بیتی صحیح نیست. شاخه کنترل پارامترها پیش از ساخت ورودی‌های حافظه، پیکربندی نامعتبر را در زمان صفر با پیام مشخص رد می‌کند. کوچک‌ترین پیکربندی معتبر یک ثبات یک‌بیتی است.

\section{روش آزمون}

هفت آزمون رفتاری مستقل نوشته شد. هر آزمون هر دو پیمانه‌ی حلقه‌دار و بدون حلقه را با
ورودی یکسان اجرا می‌کند و خروجی هرکدام را جداگانه با مقدار مورد انتظار می‌سنجد. صرف
برابرشدن دو خروجی کافی نیست؛ هر دو باید با پاسخ تعیین‌شده در آزمون برابر باشند.
ساعت آزمون‌های رفتاری دوره‌ی ۱۰ نانوثانیه دارد؛ در صورت اختلاف، شبیه‌سازی با
\کد{fatal} متوقف می‌شود و در صورت موفقیت \کد{PASS} چاپ می‌شود. هر آزمون رفتاری
فایل \لر{VCD}\LTRfootnote{Value Change Dump} تولید می‌کند.

آزمون پارامتر نامعتبر جداگانه اجرا می‌شود و معیار موفقیت آن توقف در زمان صفر با
پیام کنترل ابعاد و کد خروج غیرصفر است. در نام پیکربندی‌ها، عدد نخست تعداد ثبات‌ها
و عدد دوم عرض هر ثبات را نشان می‌دهد.

\subsection{آزمون الگوهای \لر{Wildcard} صورت آزمایش}

در \کد{tb\_tcam\_wildcard.v} برای بازسازی مثال صورت آزمایش، عرض ۸ و عمق ۴ انتخاب
شد. سه ورودی نخست در لبه‌ی ۱۵ نانوثانیه با الگوهای زیر هم‌زمان برنامه‌ریزی شدند:

\شروع{فقرات}
\فقره \لر{0110XXXX}؛
\فقره \لر{X1101XXX}؛
\فقره \لر{0X1X11X0}.
\پایان{فقرات}

الگوها با ترکیب \کد{stored\_data} و \کد{stored\_x\_mask} نمایش داده شدند. برای کلید
\لر{01101110}، همه‌ی بیت‌های محدودشده در هر سه الگو برابر بودند؛ در نتیجه سه خط نخست
فعال و \کد{match\_lines} برابر \لر{0111} شد. مجوز ورودی چهارم صفر بود؛ حتی با
وجود ماسک ورودی کاملاً بی‌تفاوت، این ورودی برنامه‌ریزی نشد و به دلیل \لر{valid} صفر
تطبیق نداشت. سپس کلید \لر{11111111} اعمال شد و هیچ الگو با آن سازگار نبود؛
بنابراین خروجی به \لر{0000} تغییر کرد.

خطوط ۵۰ تا ۵۹ پرونده‌ی \کد{tb\_tcam\_wildcard.v} بردارهای داده و ماسک مثال صورت
آزمایش و دو نتیجه‌ی مورد انتظار را پیاده‌سازی می‌کنند.

\begin{dsdlcodeblock}
\begin{verbatim}
50  // Rows 0..2: 0110XXXX, X1101XXX, 0X1X11X0.
51  // Row 3 stays invalid even though its input mask is all-X.
52  write_data = {8'b00000000, 8'b00101100, 8'b01101000, 8'b01100000};
53  write_x_mask = {8'b11111111, 8'b01010010, 8'b10000111, 8'b00001111};
54  write_enable = 4'b0111;
55  @(posedge clk);
56  #1 write_enable = 4'b0000;
57
58  check_search(8'b01101110, 4'b0111);
59  check_search(8'b11111111, 4'b0000);
\end{verbatim}
\end{dsdlcodeblock}

\شروع{شکل}[H]
\centerimg{wildcard-example.png}{14cm}
\شرح{تطبیق هم‌زمان کلید \لر{01101110} با سه الگوی سه‌حالته‌ی صورت آزمایش}
\برچسب{شکل:وایلدکارد}
\پایان{شکل}

در \کد{tb\_tcam\_wildcard.vcd}، سه داده و ماسک متفاوت در لبه‌ی ۱۵ نانوثانیه
ثبت می‌شوند. کلید \لر{01101110} از زمان ۱۶ نانوثانیه در هر دو پیاده‌سازی
خروجی \لر{0111} ایجاد می‌کند. کلید \لر{11111111} از زمان ۲۶ نانوثانیه خروجی
را به \لر{0000} تغییر می‌دهد؛ ورودی چهارم در تمام آزمون نامعتبر می‌ماند.
بنابراین بیت \لر{X} از مقایسه حذف می‌شود، اما بیت‌های صفر و یک ذخیره‌شده
باید دقیقاً با کلید برابر باشند.

\subsection{ثبت شانزده مقدار متفاوت در یک لبه ساعت}

آزمون \کد{tb\_tcam\_write.v} با پیکربندی ۱۶×۱۶، مقادیر مبنای شانزده \کد{0001} تا \کد{0010} را به‌ترتیب در ورودی‌های صفر تا ۱۵ قرار داد. تمام مجوزها در زمان ۱۰ نانوثانیه فعال شدند و همه داده‌ها در لبه ۱۵ نانوثانیه ثبت شدند. پیش از این لبه هیچ تطبیقی مجاز نبود. سپس هر شانزده مقدار جداگانه جست‌وجو شد و فقط بیت متناظر فعال شد؛ جست‌وجوی صفر نیز هیچ تطبیقی نداشت.

\شروع{شکل}[H]
\centerimg{write-16x16.png}{14cm}
\شرح{ثبت هم‌زمان شانزده داده متفاوت و جست‌وجوی مستقل آن‌ها}
\برچسب{شکل:نوشتن‌هم‌زمان}
\پایان{شکل}

در این شکل از \کد{tb\_tcam\_write.vcd}، داده‌های متفاوت ورودی‌های صفر، یک، هفت و پانزده در یک لبه ثبت شده‌اند. نتایج \کد{match\_loop} و \کد{match\_no\_loop} برای همه کلیدها یکسان‌اند و از \کد{0001} تا \کد{8000} پیش می‌روند؛ در پایان هر دو صفر می‌شوند.

\subsection{شانزده ثبات شانزده‌بیتی با ماسک‌های مستقل و نوشتن انتخابی}

آزمون \کد{tb\_tcam\_16x16.v} با عرض ۱۶ و تعداد ۱۶، هر دو پیاده‌سازی را در ۷۲ جست‌وجو بررسی کرد. در ورودی‌های صفر تا ۱۵، مقادیر مبنای شانزده \کد{0000}، \کد{1111}، \کد{2222} و به همین ترتیب تا \کد{FFFF} قرار گرفتند. ماسک ورودی $i$ فقط بیت $i$ را یک می‌کند؛ بنابراین ماسک‌ها به‌ترتیب \کد{0001}، \کد{0002}، \کد{0004} و تا \کد{8000} هستند. داده‌ها و ماسک‌ها در یک لبه، در زمان ۱۵ نانوثانیه ثبت شدند. پیش از این لبه، هیچ ورودی نباید تطبیق تولید کند.

خطوط ۶۵ تا ۷۴ پرونده‌ی \کد{tb\_tcam\_16x16.v} داده‌ها و ماسک‌های مستقل همه‌ی ورودی‌ها را مشخص می‌کنند:

\begin{dsdlcodeblock}
\begin{verbatim}
65  // Row 0 is the rightmost word. Each row ignores a different bit.
66  write_data = {16'hFFFF, 16'hEEEE, 16'hDDDD, 16'hCCCC,
67                16'hBBBB, 16'hAAAA, 16'h9999, 16'h8888,
68                16'h7777, 16'h6666, 16'h5555, 16'h4444,
69                16'h3333, 16'h2222, 16'h1111, 16'h0000};
70  write_x_mask = {16'h8000, 16'h4000, 16'h2000, 16'h1000,
71                  16'h0800, 16'h0400, 16'h0200, 16'h0100,
72                  16'h0080, 16'h0040, 16'h0020, 16'h0010,
73                  16'h0008, 16'h0004, 16'h0002, 16'h0001};
74  write_enable = 16'hFFFF;
\end{verbatim}
\end{dsdlcodeblock}

برای هر ورودی سه کلید ثابت بررسی شد: داده اصلی، داده با تغییر بیت ماسک‌شده و داده با تغییر یک بیت غیرماسک‌شده. دو کلید نخست فقط خط همان ورودی را فعال می‌کنند، اما کلید سوم هیچ تطبیقی ندارد. این بخش ۴۸ جست‌وجو دارد و همه موقعیت‌های ماسک، از بیت صفر تا بیت ۱۵، را پوشش می‌دهد. مقادیر مورد انتظار به‌صورت ثابت در آزمون نوشته شده‌اند و از منطق مقایسه پیمانه محاسبه نمی‌شوند. همه مقدارهای داده، ماسک و نتیجه در تصاویر این بخش بر حسب مبنای شانزده هستند.

\شروع{شکل}[H]
\centerimg{16x16-rows-0-3.png}{14cm}
\شرح{ماسک‌های مستقل ورودی‌های صفر تا سه در پیکربندی شانزده در شانزده}
\برچسب{شکل:شانزده‌صفر‌سه}
\پایان{شکل}

نمای ورودی‌های صفر تا سه از \کد{tb\_tcam\_16x16.vcd}، ثبت هم‌زمان چهار جفت داده و ماسک را در لبه ۱۵ نانوثانیه نشان می‌دهد. برای نمونه، \کد{0000} و \کد{0001} هر دو خروجی \کد{0001} می‌دهند، اما \کد{0002} خروجی صفر دارد. داده هر ورودی بلافاصله بالای ماسک همان ورودی قرار دارد و هر دو نسخه همین ترتیب را دارند.

\شروع{شکل}[H]
\centerimg{16x16-rows-4-7.png}{14cm}
\شرح{تطبیق دقیق، بی‌تفاوت و عدم تطبیق ورودی‌های چهار تا هفت}
\برچسب{شکل:شانزده‌چهار‌هفت}
\پایان{شکل}

در بازه ۱۳۶ تا ۲۵۶ نانوثانیه، ورودی‌های چهار تا هفت با ماسک‌های \کد{0010}، \کد{0020}، \کد{0040} و \کد{0080} بررسی می‌شوند. برای ورودی هفت، کلیدهای \کد{7777} و \کد{77F7} خروجی \کد{0080} دارند و \کد{7677} تطبیق ندارد. ثابت‌ماندن داده و ماسک هنگام تغییر کلید، ترکیبی‌بودن جست‌وجو را نشان می‌دهد.

\شروع{شکل}[H]
\centerimg{16x16-rows-8-11.png}{14cm}
\شرح{بررسی ماسک‌های نیمه بالایی کلمه در ورودی‌های هشت تا یازده}
\برچسب{شکل:شانزده‌هشت‌یازده}
\پایان{شکل}

این نما در بازه ۲۵۶ تا ۳۷۶ نانوثانیه، ماسک‌های \کد{0100}، \کد{0200}، \کد{0400} و \کد{0800} را برای ورودی‌های هشت تا یازده نمایش می‌دهد. برای نمونه، \کد{8888} و \کد{8988} هر دو خط هشت را فعال می‌کنند، ولی \کد{8A88} خروجی صفر می‌دهد. بنابراین بیت‌های ماسک در نیمه بالایی کلمه نیز به‌درستی عمل می‌کنند.

\شروع{شکل}[H]
\centerimg{16x16-rows-12-15.png}{14cm}
\شرح{بررسی ورودی‌های دوازده تا پانزده و پرارزش‌ترین بیت ماسک}
\برچسب{شکل:شانزده‌دوازده‌پانزده}
\پایان{شکل}

در بازه ۳۷۶ تا ۴۹۶ نانوثانیه، داده و ماسک ورودی‌های دوازده تا پانزده به‌ترتیب نمایش داده شده‌اند. ماسک آخرین ورودی \کد{8000} است؛ در نتیجه \کد{FFFF} و \کد{7FFF} هر دو خروجی \کد{8000} دارند، اما \کد{FFFE} تطبیق ندارد. این بخش، آخرین ورودی و پرارزش‌ترین بیت ماسک را در هر دو نسخه بررسی می‌کند.

در مرحله دوم، ورودی‌های داده و ماسک همه ثبات‌ها در زمان ۵۰۰ نانوثانیه تغییر کردند، اما \کد{write\_enable} برابر \کد{8102} بود؛ یعنی فقط ورودی‌های یک، هشت و پانزده مجوز نوشتن داشتند. در لبه ۵۰۵ نانوثانیه، ورودی یک داده \کد{A50A} و ماسک \کد{00F0}، ورودی هشت داده \کد{A50A} و ماسک \کد{0F00}، و ورودی پانزده داده \کد{0000} و ماسک \کد{FFFF} را ذخیره کردند. ماسک آخرین ورودی تمام بیت‌ها را بی‌تفاوت می‌کند. داده و ماسک سیزده ورودی غیرفعال باید دست‌نخورده بمانند.

نتایج اصلی این مرحله چنین بودند:

\شروع{فقرات}
\فقره \کد{A50A}: خروجی \کد{8102}، تطبیق هم‌زمان ورودی‌های یک، هشت و پانزده؛
\فقره \کد{A5FA}: خروجی \کد{8002}، تطبیق ورودی‌های یک و پانزده؛
\فقره \کد{AF0A}: خروجی \کد{8100}، تطبیق ورودی‌های هشت و پانزده؛
\فقره \کد{A5FB}: خروجی \کد{8000}، تطبیق فقط با ورودی کاملاً بی‌تفاوت.
\پایان{فقرات}

\شروع{شکل}[H]
\centerimg{16x16-selective-writes.png}{14cm}
\شرح{نوشتن انتخابی در ورودی‌های یک، هشت و پانزده و حفظ ورودی‌های غیرفعال}
\برچسب{شکل:شانزده‌نوشتن‌انتخابی}
\پایان{شکل}

این نما از \کد{tb\_tcam\_16x16.vcd}، مجوز \کد{8102} و تغییر داده و ماسک ورودی‌های یک، هشت و پانزده در لبه ۵۰۵ نانوثانیه را نشان می‌دهد. جفت داده و ماسک ورودی‌های صفر و هفت نیز برای مقایسه نمایش داده شده‌اند و ثابت می‌مانند. نتایج دو پیاده‌سازی، تطبیق هم‌زمان چند ورودی و تفاوت اثر دو ماسک \کد{00F0} و \کد{0F00} را نشان می‌دهند.

پس از نوشتن انتخابی، هفت جست‌وجو برای ورودی‌های بازنویسی‌شده و سیزده جست‌وجو برای حفظ ورودی‌های غیرفعال انجام شد. برای نمونه، \کد{0001} خروجی \کد{8001} و \کد{77F7} خروجی \کد{8080} داد؛ پس ورودی‌های صفر و هفت، با وجود تغییر ورودی خارجی خود، داده و ماسک قبلی را حفظ کردند. در پایان، همه داده‌ها و ماسک‌های ورودی با مجوز نوشتن صفر تغییر کردند و چهار جست‌وجوی دیگر پس از لبه ۷۱۵ نانوثانیه همچنان نتایج قبلی را دادند. مجموع ۴۸، ۷، ۱۳ و ۴ جست‌وجو برابر ۷۲ است؛ آزمون در زمان ۷۵۶ نانوثانیه با پیام \کد{PASS} پایان یافت.

\subsection{چهار ثبات سه‌بیتی با ماسک‌های مستقل}

آزمون \کد{tb\_tcam\_4x3.v} از عرض ۳ و عمق ۴ استفاده کرد. ورودی‌های صفر تا سه به‌ترتیب الگوهای \کد{001}، \کد{X10}، \کد{10X} و \کد{111} را در یک لبه دریافت کردند. هر هشت کلید ممکن سه‌بیتی بررسی شد. برای کلیدهای \کد{000}، \کد{001}، \کد{010}، \کد{011}، \کد{100}، \کد{101}، \کد{110} و \کد{111}، بردارهای مورد انتظار به‌ترتیب \کد{0000}، \کد{0001}، \کد{0010}، \کد{0000}، \کد{0100}، \کد{0100}، \کد{0010} و \کد{1000} بودند. هر دو نسخه تمام این نتایج را تولید کردند.

\شروع{شکل}[H]
\centerimg{4x3-wildcards.png}{14cm}
\شرح{پیکربندی چهار ثبات سه‌بیتی و بررسی همه کلیدهای ممکن}
\برچسب{شکل:چهارثبات‌سه‌بیتی}
\پایان{شکل}

شکل از \کد{tb\_tcam\_4x3.vcd} ثبت هم‌زمان چهار داده و ماسک متفاوت ورودی‌های یک و دو را نشان می‌دهد. تغییر کلید بدون نوشتن دوباره، همه حالت‌های جست‌وجوی سه‌بیتی را طی می‌کند و خروجی دو پیاده‌سازی در هر مرحله با مقدار مورد انتظار برابر است.

\subsection{نوشتن انتخابی و حفظ ورودی‌های غیرفعال}

در \کد{tb\_tcam\_selective.v} ابتدا داده‌های \کد{001}، \کد{010}، \کد{011} و \کد{100} در چهار ورودی نوشته شدند. سپس با مجوز \کد{1010}، در یک لبه ورودی یک به \کد{101} و ورودی سه به \کد{XXX} تغییر کرد. داده و ماسک ورودی‌های صفر و دو با وجود تغییر ورودی خارجی آن‌ها ثابت ماندند. جست‌وجوی \کد{001} نتیجه \کد{1001}، جست‌وجوی \کد{011} نتیجه \کد{1100} و جست‌وجوی \کد{101} نتیجه \کد{1010} داد. مقادیر قدیمی \کد{010} و \کد{100} فقط با ورودی کاملاً بی‌تفاوت تطبیق داشتند. تغییر همه داده‌های نوشتن با مجوز صفر نیز هیچ ورودی را بازنویسی نکرد.

\شروع{شکل}[H]
\centerimg{selective-writes.png}{14cm}
\شرح{به‌روزرسانی هم‌زمان ورودی‌های یک و سه و حفظ ورودی‌های صفر و دو}
\برچسب{شکل:نوشتن‌انتخابی}
\پایان{شکل}

در \کد{tb\_tcam\_selective.vcd}، لبه ۲۵ نانوثانیه فقط داده ورودی‌های یک و سه را تغییر می‌دهد. ماسک ورودی سه \کد{111} می‌شود، اما داده و ماسک ورودی‌های صفر و دو ثابت می‌مانند. چندتطبیقی‌بودن خروجی و حفظ داده هنگام صفرشدن تمام مجوزها در هر دو نسخه دیده می‌شود.

\subsection{کوچک‌ترین اندازه معتبر و اولویت بازنشانی}

آزمون \کد{tb\_tcam\_1x1.v} یک ثبات یک‌بیتی را بررسی کرد. ورودی نامعتبر ابتدا با هیچ‌یک از دو کلید تطبیق نداشت. پس از ثبت یک با ماسک صفر در زمان ۳۵ نانوثانیه، فقط کلید یک تطبیق داشت. با ثبت ماسک یک در زمان ۶۵ نانوثانیه، هر دو کلید تطبیق پیدا کردند. بازنشانی در زمان ۹۰ نانوثانیه هم‌زمان با مجوز نوشتن فعال شد؛ خروجی تا لبه ۹۵ نانوثانیه معتبر ماند و سپس برای هر دو کلید صفر شد. بنابراین بازنشانی هم‌زمان است و نسبت به نوشتن اولویت دارد.

\شروع{شکل}[H]
\centerimg{1x1-reset.png}{14cm}
\شرح{آزمون یک ثبات یک‌بیتی، تطبیق بی‌تفاوت و اولویت بازنشانی}
\برچسب{شکل:یک‌ثبات‌یک‌بیتی}
\پایان{شکل}

فایل \کد{tb\_tcam\_1x1.vcd} گذار از ورودی نامعتبر به تطبیق دقیق، سپس الگوی \کد{X} و در پایان بی‌اعتبارشدن در لبه بازنشانی را نشان می‌دهد. رفتار \کد{valid} و نتیجه جست‌وجو در هر دو نسخه یکسان است.

\subsection{تعداد ورودی غیرتوان دو}

آزمون \کد{tb\_tcam\_3x5.v} سه ثبات پنج‌بیتی با داده‌های \کد{00001}، \کد{10101} و \کد{11111} ساخت. این داده‌ها هم‌زمان نوشته شدند و جست‌وجوی آن‌ها به‌ترتیب خروجی‌های \کد{001}، \کد{010} و \کد{100} داد. کلید \کد{00000} هیچ تطبیقی نداشت. این آزمون نشان می‌دهد تعداد ورودی‌ها لازم نیست توان دو باشد و مرزبندی داده‌ها برای عرض پنج بیت نیز صحیح است.

\شروع{شکل}[H]
\centerimg{3x5-depth.png}{14cm}
\شرح{نوشتن و جست‌وجوی سه ثبات پنج‌بیتی}
\برچسب{شکل:سه‌ثبات‌پنج‌بیتی}
\پایان{شکل}

فایل \کد{tb\_tcam\_3x5.vcd} ثبت سه کلمه پنج‌بیتی در لبه ۱۵ نانوثانیه و فعال‌شدن مستقل سه خط تطبیق را در هر دو نسخه ثبت می‌کند.

\subsection{رد یک ثبات صفر‌بیتی و ابعاد نامعتبر}

آزمون \کد{tb\_tcam\_invalid.v} در حالت پیش‌فرض عرض صفر و تعداد یک دارد. با \کد{NO\_LOOP=0} نسخه حلقه‌دار و با \کد{NO\_LOOP=1} نسخه بدون حلقه جداگانه بررسی شدند. هر دو در زمان صفر با کد خروج یک و پیام \کد{TCAM requires DATA\_WIDTH >= 1 and DEPTH >= 1} متوقف شدند. حالت‌های عرض منفی، تعداد صفر و تعداد منفی نیز برای هر دو پیاده‌سازی رد شدند.

اگر پارامتر نامعتبر پذیرفته شود، خود آزمون در زمان یک نانوثانیه با پیام متفاوتی خطا می‌دهد؛ بنابراین صرف دریافت هر خطایی به معنای موفقیت نیست. این آزمون شکل موج رفتاری تولید نمی‌کند، زیرا پیکربندی نامعتبر پیش از شروع کار حافظه رد می‌شود.

\شروع{شکل}[H]
\centerimg{zero-width-rejected.png}{14cm}
\شرح{خروجی واقعی رد یک ثبات صفر‌بیتی در هر دو پیاده‌سازی}
\برچسب{شکل:رد‌عرض‌صفر}
\پایان{شکل}

تصویر پایانه نشان می‌دهد هر دو پیاده‌سازی با پیام کنترل ابعاد در زمان صفر و کد خروج یک متوقف می‌شوند. فایل‌های \کد{invalid-0x1-0.log} و \کد{invalid-0x1-1.log} متن همین خطاها را نگه می‌دارند. این تصویر مدرک رد پارامتر است و شکل موج یک حافظه معتبر نیست.

\section{نتایج اجرای خودکار}

هر هفت آزمون رفتاری با \لر{Icarus Verilog} و گزینه‌های \کد{-g2012} و \کد{-Wall}
کامپایل و اجرا شدند. هیچ هشدار کامپایل یا شکست آزمونی ثبت نشد. نتیجه‌ی همه‌ی آزمون‌ها
\لر{PASS} بود:

\شروع{فقرات}
\فقره \کد{tb\_tcam\_write.v}: شانزده داده‌ی مستقل در پیکربندی $16\times16$؛
\فقره \کد{tb\_tcam\_16x16.v}: ۷۲ جست‌وجو با ماسک‌های مستقل شانزده‌بیتی و نوشتن انتخابی؛
\فقره \کد{tb\_tcam\_4x3.v}: ماسک مستقل و تمام کلیدهای پیکربندی $4\times3$؛
\فقره \کد{tb\_tcam\_selective.v}: نوشتن انتخابی و حفظ ورودی‌های غیرفعال؛
\فقره \کد{tb\_tcam\_1x1.v}: حداقل اندازه‌ی معتبر، تطبیق دقیق و \لر{X}، و اولویت بازنشانی؛
\فقره \کد{tb\_tcam\_3x5.v}: تعداد غیرتوان دو و برش صحیح داده‌ی پنج‌بیتی؛
\فقره \کد{tb\_tcam\_wildcard.v}: سه الگوی صورت آزمایش در پیکربندی $4\times8$.
\پایان{فقرات}

در هر آزمون، خروجی هر دو نسخه با مقدار مورد انتظار سنجیده شد. فایل‌های \لر{VCD}،
نمای ذخیره‌شده‌ی \لر{GTKWave} با پسوند \کد{gtkw} و گزارش‌های اجرا با همان نام آزمون
در \کد{build} قرار دارند. هشت اجرای پارامتر نامعتبر، شامل چهار حالت برای هر نسخه،
همگی پیام کنترل ابعاد را در زمان صفر تولید کردند. شکل‌های موج مستقیماً از پنجره‌ی
\لر{GTKWave} ثبت شده‌اند و هیچ نشانگر زمانی روی آن‌ها قرار ندارد. تصویر رد عرض صفر
از خروجی واقعی پایانه تهیه شده است.

\section{بررسی سنتزپذیری}

هر دو پیمانه‌ی \کد{tcam} و \کد{tcam\_no\_loop} در پنج پیکربندی $16\times16$،
$4\times3$، $1\times1$، $3\times5$ و $4\times8$ با \لر{Yosys} و گذر
\کد{synth} و سپس \کد{check -assert} بررسی شدند. هر ده اجرا بدون مشکل پایان یافت.
در پیکربندی پیش‌فرض، هر سطح بالا ۱۶ نمونه‌ی \کد{tcam\_entry} دارد.
آمار پس از \کد{synth} برای هر دو نسخه ۱۳۱۲ سلول بود؛ از جمله ۵۱۲ سلول ثبات داده و
ماسک و ۱۶ سلول ثبات اعتبار. این آمار متعلق به سلول‌های عمومی \لر{Yosys} پس از
نگاشت منطقی است و شمار \لر{LUT}\LTRfootnote{Look-Up Table} یک
\لر{FPGA}\LTRfootnote{Field-Programmable Gate Array} مشخص نیست.

\section{تصمیم‌های طراحی و ملاحظات پیاده‌سازی روی \لر{FPGA}}

نمایش سه‌حالته با دو بردار داده و ماسک انتخاب شد، زیرا استفاده از مقدار \لر{x} زبان
\لر{Verilog} به‌عنوان داده‌ی ذخیره‌شده برای سخت‌افزار واقعی مناسب نیست. در این نمایش، هر
بیت \لر{TCAM} عملاً به یک بیت داده و یک بیت ماسک نیاز دارد. برای ساختار پیش‌فرض، ۲۵۶ بیت
داده، ۲۵۶ بیت ماسک و ۱۶ بیت اعتبار لازم است.

جست‌وجو در همه‌ی ۱۶ ورودی موازی است. این ساختار پاسخ ترکیبی سریع ایجاد می‌کند، اما با افزایش
\لر{DEPTH} و \کد{DATA\_WIDTH} تعداد مقایسه‌کننده‌ها و مسیرهای منطقی افزایش می‌یابد.
از آنجا که بسیاری از \لر{FPGA}ها بلوک \لر{TCAM} اختصاصی ندارند، چنین طرحی معمولاً با
ثبات‌ها و \لر{LUT}ها پیاده‌سازی می‌شود. انتخاب بردار \کد{match\_lines} به‌جای
\لر{priority encoder} باعث حفظ همه‌ی تطبیق‌ها و ساده‌ماندن هسته‌ی \لر{TCAM} شد؛ در
صورت نیاز، اولویت‌گذاری می‌تواند به‌عنوان پیمانه‌ای جداگانه پس از \کد{match\_lines} افزوده شود.

\لر{reset} فقط \لر{valid} را پاک می‌کند. این تصمیم برای بورد مناسب است، زیرا
بی‌اعتبارکردن ۱۶ ورودی به پاک‌سازی هم‌زمان تمام ۵۱۲ بیت داده و ماسک نیاز ندارد. همچنین
بیت مجوز هر ورودی از تغییر ناخواسته‌ی داده و ماسک ورودی‌های غیرفعال جلوگیری می‌کند.

در هر دو نسخه، برای شانزده ورودی شانزده‌بیتی، ۲۵۶ خط داده، ۲۵۶ خط
ماسک و ۱۶ مجوز مستقل به هسته متصل می‌شود. این رابط امکان بارگذاری همه ورودی‌ها در
یک لبه را فراهم می‌کند و به مسیر مستقل داده، ماسک و مجوز برای هر ورودی نیاز دارد.
تعداد بیت‌های ذخیره‌سازی و منطق تطبیق در دو پیاده‌سازی یکسان است؛ بازنشانی همچنان نسبت به
همه نوشتن‌ها اولویت دارد.

\section{نتیجه‌گیری}

یک \لر{TCAM} پارامتری با مقدار پیش‌فرض ۱۶ ورودی ۱۶‌بیتی طراحی شد. هر ورودی داده، ماسک
و اعتبار مستقل دارد و جست‌وجو به‌صورت ترکیبی در تمام ورودی‌های معتبر انجام می‌شود. رابطه‌ی
منطقی استفاده‌شده، بیت‌های ماسک‌شده را حذف و فقط اختلاف بیت‌های اجباری را بررسی می‌کند.

هفت آزمون رفتاری، نوشتن هم‌زمان داده‌های متفاوت، ماسک‌های مستقل، حفظ ورودی‌های غیرفعال،
بازنشانی و اندازه‌های متفاوت را برای هر دو نسخه تأیید کردند. مثال \لر{01101110}
با هر سه الگوی \لر{0110XXXX}، \لر{X1101XXX} و \لر{0X1X11X0} تطبیق یافت
و خروجی چندتطبیقی \لر{0111} تولید شد. آزمون یک ثبات صفر‌بیتی و دیگر ابعاد نامعتبر
نیز با رد صریح پیکربندی پایان یافت. بررسی \لر{Yosys} سنتزپذیری هر دو ساختار را
در پنج اندازه‌ی آزموده‌شده تأیید کرد.
